% ==================================================
%  Capítulo 1 — Introducción
% ==================================================
\section{Introducción}\label{sec:introduccion}

Las metalentes son dispositivos ópticos ultradelgados que permiten enfocar la luz sin recurrir a la curvatura de los lentes convencionales. Su principio de operación se basa en la manipulación de la fase de la luz a escala nanométrica mediante arreglos de nanoestructuras distribuidas sobre una superficie plana. Un ejemplo relevante lo constituyen las metalentes formadas por nano-rendijas, ya sean metálicas o dieléctricas, donde cada rendija introduce un retraso de fase específico sobre la onda incidente. Al distribuir estratégicamente estas nano-rendijas a lo largo de la superficie, es posible construir un perfil de fase capaz de concentrar o redirigir la luz hacia una región de interés.

Las metalentes plasmónicas~\cite{khorasaninejad2017metalenses}, fabricadas típicamente con metales nobles como el oro o la plata depositados sobre un sustrato dieléctrico, aprovechan resonancias de plasmones superficiales, asociadas a oscilaciones colectivas de los electrones libres del metal, para confinar y redirigir la luz a escala nanométrica. Predecir la manera en que la fase y la amplitud de la luz son modificadas por estas nanoestructuras resulta fundamental para el diseño de dispositivos fotónicos eficientes. Este problema, junto con el alto costo computacional de su simulación electromagnética directa, motiva el presente trabajo.

\subsection{Antecedentes}\label{subsec:antecedentes}

Las metasuperficies, entendidas como arreglos planares de nanoestructuras sub-longitud de onda, han emergido como una alternativa versátil a la óptica refractiva convencional, al permitir controlar de manera flexible la fase, la amplitud y la polarización de un frente de onda mediante componentes de espesor nanométrico~\cite{khorasaninejad2017metalenses}. Entre sus aplicaciones, las metalentes, es decir, metasuperficies diseñadas para imponer un perfil de fase convergente, han atraído especial interés por su potencial para reemplazar sistemas de lentes voluminosos en imagenología, sensado y fotónica integrada.

Una familia particular es la de las metalentes plasmónicas de nano-rendijas: una película metálica delgada, típicamente de oro, perforada por rendijas de ancho sub-longitud de onda. La luz atraviesa estas aberturas mediante transmisión óptica extraordinaria~\cite{ebbesen1998eot}, y cada rendija puede modelarse como una guía de onda plasmónica metal-aislante-metal que introduce un retraso de fase $\phi = \beta t$, dependiente de sus propiedades geométricas. Al distribuir adecuadamente los anchos de las rendijas a lo largo de la superficie, es posible aproximar un perfil de fase objetivo y estructurar así el frente de onda transmitido. El trabajo de referencia de este proyecto~\cite{airy_metalens} utiliza un arreglo de $101$ nano-rendijas de ancho variable en una película de oro para generar un haz de Airy en espacio libre, y establece el marco físico del que parte el presente estudio.

El obstáculo central para el diseño de estos dispositivos es el costo de la simulación electromagnética. Resolver las ecuaciones de Maxwell sobre una geometría de nano-rendijas mediante métodos numéricos como FDTD o el método de elemento finito (FEM) requiere de minutos a horas por configuración, lo que vuelve impráctica la exploración exhaustiva del espacio de diseño y favorece, históricamente, los enfoques guiados por intuición física o por barridos paramétricos limitados. En este contexto, el diseño inverso, entendido como la búsqueda de una geometría que produzca una respuesta óptica deseada, ha cobrado fuerza como paradigma para abordar problemas de optimización en nanofotónica~\cite{molesky2018inverse}.

El aprendizaje automático ofrece dos vías complementarias para apoyar este tipo de tareas. La primera son las redes neuronales para diseño inverso directo, que buscan aprender el mapeo de una respuesta deseada hacia una geometría. Un reto conocido de este enfoque es la no-unicidad del problema inverso, que arquitecturas en tándem, al acoplar un modelo directo y uno inverso, han logrado mitigar parcialmente~\cite{liu2018tandem}. La segunda vía son los modelos subrogados (\textit{surrogate models}), que aproximan la respuesta de un simulador costoso a una fracción de su tiempo de evaluación y se acoplan a un optimizador para explorar regiones prometedoras del espacio de diseño antes de recurrir nuevamente al simulador de referencia. Esta combinación ha mostrado resultados relevantes tanto en la predicción de parámetros de nanoestructuras plasmónicas mediante modelos de ensamble~\cite{maia2025inverse} como en el diseño de metasuperficies acoplando subrogados con algoritmos evolutivos~\cite{gao2025surrogate}.

El presente trabajo se inscribe en esta segunda línea. Se construye un modelo subrogado del campo óptico focal de metalentes plasmónicas de nano-rendijas y se acopla a un algoritmo genético para guiar la búsqueda de geometrías candidatas, reservando la simulación FDTD para validar y refinar los candidatos más prometedores.

\subsection{Planteamiento del Problema}\label{subsec:planteamiento}

Formalmente, el problema que aborda este trabajo es el diseño inverso asistido de una metalente plasmónica de nano-rendijas. Dada una respuesta óptica deseada, en este caso un perfil de intensidad que concentre la luz en una región focal para $\lambda = 630$~nm, se busca identificar configuraciones geométricas de rendijas capaces de aproximar dicha respuesta. En particular, el diseño queda parametrizado por un conjunto de anchos de rendija y una separación centro a centro, cuya combinación determina la fase y la amplitud del campo transmitido.

Este problema presenta dos dificultades principales. La primera es el alto costo del evaluador directo, ya que cada evaluación de una geometría requiere una simulación electromagnética FDTD, por lo que una búsqueda exhaustiva o la evaluación directa de todos los candidatos generados por un optimizador poblacional resulta impráctica dentro de los recursos computacionales disponibles. La segunda dificultad es la no linealidad del mapeo geometría--respuesta óptica y la baja densidad de soluciones de alto desempeño. En el espacio de diseño considerado, las configuraciones que producen un enfoque satisfactorio representan una fracción muy pequeña de las geometrías posibles, de modo que un muestreo no dirigido rara vez encuentra candidatos cercanos al criterio deseado.

Frente a estas dificultades, el presente proyecto no busca reemplazar de manera global la simulación FDTD, sino utilizar un modelo subrogado como aproximador rápido durante la etapa de búsqueda. El surrogate permite evaluar y priorizar un gran número de geometrías candidatas a bajo costo, mientras que la simulación FDTD se conserva como referencia física para validar y refinar los candidatos más prometedores. Por ello, el enfoque se plantea como un pipeline heurístico de búsqueda local asistida por aprendizaje automático, especialmente útil para explorar regiones cercanas a geometrías previamente observadas o validadas.

El criterio operativo de éxito, heredado del trabajo de referencia y verificado mediante FDTD, se define como un ajuste parabólico al perfil de intensidad focal con coeficiente de determinación $R^2 \geq 0.90$. Este criterio resume cuantitativamente la regularidad del perfil de intensidad asociado al enfoque en el marco de este estudio. No obstante, se reconoce que $R^2$ no agota la caracterización física de una metalente, por lo que debe interpretarse como una métrica práctica de selección y comparación de candidatos, complementada con la validación electromagnética directa y el análisis de estabilidad local.

\subsection{Justificación}\label{subsec:justificacion}

La pertinencia de este proyecto descansa en tres argumentos. En primer lugar, un argumento de eficiencia computacional: dado que cada simulación FDTD de una metalente cuesta del orden de decenas de minutos de cómputo, optimizar el diseño mediante evaluación directa del simulador resulta inviable para un algoritmo poblacional, pues una sola corrida requeriría evaluar cientos o miles de geometrías. Un modelo subrogado que aproxima la respuesta óptica en tiempos del orden de milisegundos permite reducir el costo de exploración en varios órdenes de magnitud y concentrar el presupuesto de simulación FDTD en un número reducido de candidatos seleccionados.

En segundo lugar, un argumento metodológico y científico: la combinación de alta dimensionalidad geométrica y escasez de soluciones de alto desempeño convierte a este problema en un caso de estudio exigente para evaluar hasta qué punto un pipeline de aprendizaje automático puede guiar la búsqueda hacia regiones prometedoras, aun cuando el modelo subrogado no sea suficientemente preciso para sustituir al simulador en todo el espacio. Más allá de alcanzar un umbral puntual de calidad, interesa distinguir entre soluciones de alto desempeño pero frágiles y regiones de diseño con mayor robustez relativa frente a perturbaciones geométricas.

Finalmente, un argumento de transferencia de conocimiento: los productos del trabajo, un conjunto de datos depurado, un modelo subrogado entrenado y un repositorio que integra el flujo de preprocesamiento, modelado, búsqueda y validación, constituyen una base inicial reutilizable para la línea de investigación en diseño computacional de dispositivos fotónicos. Esta base puede servir para extender el estudio hacia otros perfiles de fase, longitudes de onda, restricciones de fabricación o estrategias de optimización, siempre considerando que la generalización fuera del dominio simulado requiere validación adicional.

\subsection{Objetivos}\label{subsec:objetivos}

\subsubsection{Objetivo General}

Desarrollar y evaluar un prototipo computacional de diseño inverso asistido por aprendizaje automático y optimización para metalentes plasmónicas de nano-rendijas, orientado a acelerar la búsqueda de geometrías candidatas que produzcan enfoque óptico a $\lambda = 630$~nm, empleando el modelo subrogado para priorizar candidatos y reservando la simulación FDTD como mecanismo indispensable de validación física.

\subsubsection{Objetivos Específicos}

\begin{enumerate}
\item Construir y depurar un conjunto de datos mediante la ejecución sistemática de simulaciones FDTD, vinculando los parámetros geométricos de las nano-rendijas con descriptores de su respuesta óptica, en particular el campo focal complejo, la amplitud, la fase y métricas derivadas de calidad de enfoque.

\item Implementar y evaluar modelos de regresión bajo el enfoque de \textit{Surrogate Modeling}, analizando tanto su capacidad predictiva como sus limitaciones para aproximar la respuesta óptica en el espacio de diseño muestreado.

\item Desarrollar un algoritmo de optimización asistido por el modelo subrogado para explorar regiones prometedoras del espacio de diseño, priorizar geometrías candidatas y reducir el número de simulaciones FDTD requeridas durante la búsqueda.

\item Validar mediante simulación FDTD directa las geometrías recomendadas por el pipeline, cuantificar su calidad de enfoque mediante el criterio operativo $R^2 \geq 0.90$ y analizar su estabilidad local frente a perturbaciones geométricas de escala nanométrica.
\end{enumerate}


\subsection{Alcance y Limitaciones}\label{subsec:alcance}

El alcance del proyecto se delimita a la exploración, implementación y validación numérica de un pipeline computacional para el diseño inverso asistido de metalentes plasmónicas. El trabajo opera exclusivamente en el dominio de la simulación y se organiza en cuatro ejes: (i) generación y depuración de datos mediante simulaciones FDTD, (ii) entrenamiento y evaluación de modelos subrogados, (iii) implementación de estrategias de optimización asistidas por dichos modelos, y (iv) validación numérica de candidatos seleccionados contra el simulador electromagnético.

El trabajo no incluye procesos de fabricación física, litografía ni manufactura de dispositivos; tampoco contempla caracterización experimental en laboratorio óptico. La validación se restringe a la comparación numérica contra la respuesta calculada por FDTD, que se considera la referencia física del estudio. Asimismo, no se desarrollan solucionadores de ecuaciones de Maxwell desde cero, sino que se integran y automatizan herramientas numéricas preexistentes dentro del flujo de trabajo computacional.

Una limitación central del proyecto es el tamaño y la cobertura del conjunto de datos. Aunque se consolidaron $1{,}110$ simulaciones FDTD, el espacio de diseño es de alta dimensión y está dominado por configuraciones de bajo desempeño. Las geometrías que alcanzan valores altos de calidad focal constituyen una fracción muy reducida del dataset, lo que limita la capacidad del modelo subrogado para aprender de manera global las regiones óptimas. En consecuencia, el surrogate no debe interpretarse como un predictor confiable de geometrías óptimas en todo el espacio de diseño, sino como una herramienta útil para aproximar tendencias, priorizar candidatos y acelerar búsquedas locales en regiones cercanas a geometrías previamente observadas o validadas.

Otra limitación proviene del diseño de experimentos. Para hacer viable la generación de datos dentro del periodo de la estancia, las geometrías fueron restringidas a familias parametrizadas de perfiles de ancho, manteniendo una separación centro a centro constante por metalente. Esta decisión redujo la dimensionalidad efectiva del problema y permitió construir un dataset utilizable, pero también implica que el espacio explorado no cubre todas las posibles configuraciones de una metalente con $101$ rendijas independientes. Por tanto, los resultados obtenidos no garantizan optimalidad global ni generalización hacia diseños arbitrarios fuera del manifold muestreado.

El estudio también se limita a una configuración física específica: longitud de onda fija de $\lambda = 630$~nm, simulaciones bidimensionales, polarización TM, materiales y parámetros geométricos definidos, y arreglos de $101$ nano-rendijas. No se evalúa la generalización hacia otras longitudes de onda, otros materiales, diferentes números de rendijas, separaciones variables por par de rendijas, geometrías libres ni optimización topológica. Cualquier extensión hacia estos escenarios requeriría generar nuevos datos, reentrenar o adaptar el modelo y validar nuevamente mediante FDTD.

Finalmente, el criterio $R^2 \geq 0.90$ se utiliza como una métrica operativa para seleccionar y comparar candidatos con base en el ajuste parabólico del perfil de intensidad focal. Aunque este criterio permite cuantificar de manera práctica la calidad de enfoque dentro del marco del proyecto, no agota la caracterización óptica de una metalente. Métricas adicionales como eficiencia de transmisión, posición focal, anchura del lóbulo principal, lóbulos secundarios y tolerancia a fabricación deberían incorporarse en etapas posteriores para una evaluación más completa del desempeño físico del dispositivo.

\subsection{Productos Esperados}\label{subsec:productos}

Como resultado de la estancia se entregan: (i)~un dataset estructurado y depurado que consolida geometrías de metalentes de nano-rendijas y sus respuestas ópticas simuladas mediante FDTD; (ii)~un conjunto de modelos subrogados entrenados y evaluados, incluyendo el análisis del enfoque inicial basado en la predicción directa de $R^2$ y la posterior reformulación hacia la predicción del campo óptico focal complejo; (iii)~un pipeline computacional en Python que integra preprocesamiento de datos, entrenamiento del surrogate, optimización asistida, selección de candidatos y validación por simulación FDTD; (iv)~un conjunto de geometrías candidatas recomendadas y confirmadas numéricamente, junto con el análisis de su calidad de enfoque y estabilidad local; y (v)~el presente reporte técnico, que documenta la metodología, los resultados obtenidos, las limitaciones identificadas y las oportunidades de mejora para trabajos futuros.
